\section{Results and Discussion}

This chapter reports the repeated-task study and the completed 200-run CloudraBench comparison, then connects the observed configuration differences to the research trajectories and local diagnostic experiments.

\input{sections/03_benchmark_results}
\FloatBarrier

\subsection{From chemical intent to reliable geometric operations}

\subsubsection{Molecular identity provides a persistent reference for operations}
\begin{figure*}[tbp]
\centering
\includegraphics[width=\textwidth]{figures/fig06.pdf}
\caption{Molecular-identity judgments under changes in molecular representation and option order. Each of the 16 molecular pairs has four input conditions. Circles and diamonds indicate correct and incorrect judgments, respectively; the letters retain the model's actual answers. Shaded regions denote identical-string conditions. The right panel shows the two C\textsubscript{5}H\textsubscript{8} connectivities in pair 14 and the meanings of the answer options.}
\label{fig:06}
\end{figure*}


A geometric modification first requires identifying the molecule, atoms, and reaction center to which it applies. A molecule can be represented by different SMILES traversals, while a molecular formula can describe different connectivities. To examine how the model handles these distinctions, we constructed 16 molecular-identity pairs covering identical identities, connectivity changes within the same molecular formula, changes in connectivity and charge location, double-bond stereochemistry, and protonation changes. Each pair was presented using the original strings or equivalent rewritings and with two answer orders, yielding 64 judgments.

The model answered 58 correctly. Equivalent rewriting changed the semantic judgment in 4 of 32 matched representation comparisons; reversing answer order changed the judgment in 6 of 32 matched option-order comparisons. All six of the latter cases involved repeated selection of the letter A. One C\textsubscript{5}H\textsubscript{8} pair in Figure~\ref{fig:06} illustrates this behavior directly: an endocyclic and an exocyclic double bond define different connectivities, yet the model answered A under all four conditions, although A meant “same” in one order and “different” in the other. This result locates representation sensitivity in judgments about the chemical objects themselves and indicates the need for stable identity references throughout a research process.

Cloudra assigns this role to molecular topology, atom mapping, and reaction centers. These references allow subsequent measurements, fragment movements, and endpoint comparisons to remain attached to the originally specified objects as the model moves between descriptions and configurations. Atom correspondence also connects product-endpoint alignment, path interpolation, and reaction-center monitoring, allowing local geometric operations to proceed with a consistent set of chemical identities.



\subsubsection{Geometric measurements reduce repeated numerical reasoning}

After identifying the research objects, the model must extract spatial relationships relevant to an operation from their coordinates. We selected four actual structures from the NB1 trajectories and asked the model to compare two specified interatomic distances. Each structure was presented with both answer orders and with either a complete XYZ geometry or program-measured distances. These two information conditions were further crossed with direct-answer and thinking modes, producing 32 native chat calls (Figure~\ref{fig:07}).

In direct-answer mode, all eight responses to XYZ inputs selected A, with four correct answers. Providing measured distances produced eight correct answers and preserved the semantic answer when option order was reversed. With thinking enabled, the XYZ condition produced seven correct answers; the remaining response reached the output limit of 1,024 generated tokens. Increasing the budget for that input alone yielded a correct answer within 1,353 generated tokens. All eight answers in the measured-distance condition were correct. Across the eight paired comparisons with thinking enabled, the median generated token counts were 807.5 for XYZ inputs and 86 for measured distances.

These results suggest that geometric judgment combines atom identification, numerical calculation, and relational comparison. Additional reasoning can perform these local calculations, while program measurements directly supply the relational quantities needed for the subsequent decision. Cloudra therefore allows the model to express objectives through bond lengths, bond angles, dihedral angles, and fragment relationships, and uses deterministic operators to change the coordinates. This division of work assigns repeated numerical operations to the program and makes the resulting changes directly inspectable, consistent with the allocation of responsibilities in program-aided language models.\citep{ref15} Descriptions of steric crowding, reactant preorganization, and local fragment rotation build on this foundation to use spatial relationships in selecting and constructing candidate configurations.


\begin{figure*}[tbp]
\centering
\includegraphics[width=.64\textwidth]{figures/fig07.pdf}
\caption{Geometry judgments under coordinate inputs and measured-distance inputs. Four source structures are compared under direct response and reasoning-enabled conditions, each with original (F) and reversed (R) option order. Circles indicate correct judgments, diamonds incorrect judgments, and the triangle a response that reached the output budget. Letters retain the actual answers.}
\label{fig:07}
\end{figure*}



\begin{figure}[tbp]
\centering
\includegraphics[width=\linewidth]{figures/fig08.pdf}
\caption{Generated-token counts with reasoning enabled. Each line connects Cartesian-coordinate input and measured-distance input for the same source structure and option order; eight paired comparisons are shown. The dashed line marks the output limit.}
\label{fig:08}
\end{figure}


\subsubsection{Feedback must match the specified atoms and geometric objective}
\begin{figure*}[tbp]
\centering
\includegraphics[width=.9\textwidth]{figures/fig09.pdf}
\caption{Target-specific geometry operations in E1. Newman projections retain the directions of the original candidate coordinates viewed down C2\(\to\)C3. H7 is the specified transferring atom. Candidate A and candidate B give target angles of 59.4° and 179.4°, respectively. The operation sequence is shown alongside the projections.}
\label{fig:09}
\end{figure*}


Deterministic execution must remain aligned with chemical intent. An actual E1 trajectory provides a complete example. The model sought to bring the dihedral angle defined by the transferring hydrogen H7, C2, C3, and leaving atom Cl4 close to 180°. It rotated the fragment on the C3 side about the C2–C3 axis by +120° and \( - \)120°, producing candidates A and B, and set the C3–Cl4 distance to 2.35 Å. The representative dihedral in the default structure report involved H5, whereas the current objective involved H7. The model then explicitly measured H7–C2–C3–Cl4 in both candidates, obtaining approximately 59.4° and 179.4°, and used candidate B in its next actual submission.

Figure~\ref{fig:09} shows Newman projections along C2\(\to\)C3 drawn directly from the saved candidate coordinates, and Figure~\ref{fig:10} summarizes the corresponding target-angle residuals. Candidate B places the H7 target angle within 0.6° of 180°, compared with a residual of 120.6° for A. The corresponding residuals for the H5 comparison were 124.2° and 115.8° in A and B, respectively. These measurements reveal the object dependence of feedback: a reading acquires its chemical meaning through its correspondence to the currently specified atoms and operational objective.

Local comparisons using the same measurements further connect information to action. When only H5 measurements were provided, the model requested an additional H7 measurement under all three cyclic option permutations. Adding H7 measurements led it to select B under all three permutations. Two further paired tests of information sufficiency also shifted from requesting additional measurements to judging the information sufficient when the target quantity was added. Atom mapping, internal coordinates, fragment operations, and feedback can therefore be understood as one sequence of research actions: identify the target object, make the geometric change, measure the target quantity, and select the next candidate. Preorganization, rigid-body alignment, and path preparation use the same correspondence to carry reaction intent into subsequent exploration of the potential energy surface.




\begin{figure*}[tbp]
\centering
\begin{minipage}[t]{.48\textwidth}
\vspace{0pt}\centering
\includegraphics[width=\linewidth]{figures/fig10.pdf}
\caption{Circular angular residuals from the specified 180° target for candidates A and B. H7 is the specified target atom; H5 is the comparator.}
\label{fig:10}
\end{minipage}\hfill
\begin{minipage}[t]{.48\textwidth}
\vspace{0pt}\centering
\includegraphics[width=\linewidth]{figures/fig11.pdf}
\caption{Next-action choices under three local option permutations. When only H5 measurements are supplied, the model requests H7 measurements; with H7 supplied, it selects candidate B.}
\label{fig:11}
\end{minipage}
\end{figure*}




\subsubsection{Internal transfer of measured facts affects selection}

To examine how the model uses supplied measurements, we performed activation replacement at a single token position in the same local quantized model.\citep{ref16} One distance in the input was changed from 2.029 Å to 1.029 Å, while the other remained at 1.248 Å, reversing the correct comparison answer. We then transferred the residual-stream state at the numeric token position from the original input into the modified input at layers 7, 15, 31, and 47, and read the full-vocabulary probability assigned to the answer supported by the original fact.

The effect varied with layer and option order (Figure~\ref{fig:13}). With reversed options, the probability of A, the original-fact answer, increased from 0.074 for the modified input to 0.956 and 0.718 after replacement at layers 7 and 15, respectively. Replacements at layers 31 and 47 yielded 0.072 and 0.070. With the original option order, the probability of answer B changed from 0.009 to 0.029 and 0.015 after replacement at the first two tested layers, while the last two remained close to the modified-input baseline. Self-replacement controls and control replacements at an unchanged prefix position produced zero logit differences across the full vocabulary.

This internal intervention extends the analysis from supplying measurements to examining the causal mechanisms governing how measurement information shapes decisions: under reversed options, intermediate-layer state replacements dramatically flipped model predictions, demonstrating that localized numerical representations directly participate in downstream choices. Under the original option order, however, the patch produced only a minor shift, aligning closely with the behavioral positional bias identified in Section 3.2.1—when the target answer occupies the non-preferred option B, deep attention may be influenced by the shallow positional prior. This asymmetry suggests at the representation level that injecting ungrounded numerical facts into free-form text prompts may be vulnerable to attentional attenuation and prior biases; by coupling deterministic structural operators that explicitly bind geometric measurements, chemical entities, and operational semantics a stable scientific decision loop can be established. Cloudra's structured measurements thus serve not only as the outputs of geometric calculation, but as physical anchors for subsequent mechanistic reasoning.


\begin{figure*}[tbp]
\centering
\includegraphics[width=.9\textwidth]{figures/fig12.pdf}
\caption{Intervention on a single reported distance. A one-digit edit changes the C2–C4 distance from 2.029 to 1.029 Å while O3–H8 remains 1.248 Å. The original-input residual state at token position 53 is inserted into the changed-input run at a specified layer. Token positions are zero-based.}
\label{fig:12}
\end{figure*}



\begin{figure*}[tbp]
\centering
\includegraphics[width=.9\textwidth]{figures/fig13.pdf}
\caption{Original-fact answer probability after state replacement. Curves show full-vocabulary probability mass for the original-fact answer under forward (B) and swapped (A) option order, including bare and leading-space token spellings. Horizontal lines show original-input and changed-input baselines. Layer indices are zero-based.}
\label{fig:13}
\end{figure*}



\begin{figure*}[tbp]
\centering
\includegraphics[width=.75\textwidth]{figures/fig14.pdf}
\caption{Control substitutions for the state-intervention experiment. Same-state self replacement at p53 and replacement of the unchanged prefix at p0 preserve all vocabulary logits in both option orders. The table reports the maximum absolute logit difference. Token positions are zero-based.}
\label{fig:14}
\end{figure*}


\subsection{From local computational results to evidence-supported scientific delivery}

\subsubsection{The DAG preserves scientific support relationships among results}

Once geometric objects enter computation, optimized structures, energies, frequencies, and paths must be assembled into support for the same reaction step. The local validity and attribution of each result jointly determine this support chain: frequencies must correspond to the declared saddle point, the IRC must start from that saddle point, and optimizations and frequencies for the final frames on both sides must correspond to the respective endpoints. The DAG records these relationships on species and elementary steps, providing a persistent research index as results accumulate. Execution state, numerical convergence, and scientific support are recorded separately, allowing subsequent work to wait for a calculation, revise and rerun it, or continue assembling evidence.

We extracted four source evidence chains from the archive and constructed versions that preserved the links or changed a single reference. Crossing these with flat and linked representations and two answer orders yielded 32 inputs. The altered references concerned the frequency geometry, IRC starting geometry, endpoint-optimization input, or target species; local termination, convergence, energy, and frequency values were held constant. Figure~\ref{fig:15} illustrates the E1 case: changing the frequency geometry reference from the saddle point to an existing endpoint changes how that frequency result supports the current step.

All 16 judgments with linked representations were correct, compared with 14 of 16 with flat representations. Both errors occurred when a reference had been changed and the options reversed, in E1 and the second NB1 step. In subsequent native follow-up calls on two selected source chains, enabling thinking yielded correct judgments for all 16 judgments of the evidence relationships, including the two previously incorrect cases. This result suggests that explicit links reduce the model's work in locating and assembling computational relationships, consistent with observations that graph encoding affects relational reasoning in language models.\citep{ref17} In Cloudra, this representation is further bound to actual computational artifacts, allowing scientific facts to accumulate continuously for the same species, step, and attempt.

\begin{figure*}[tbp]
\centering
\begin{minipage}[t]{.48\textwidth}
\vspace{0pt}\centering
\includegraphics[width=\linewidth]{figures/fig15.pdf}
\caption{Evidence–geometry association in the E1 source chain. The controlled input changes the frequency geometry reference from g\_63 to g\_42 while retaining local termination, convergence, energy, and frequency values. The diagram preserves the links from TS optimization through bidirectional IRC to endpoint optimization and frequency calculations.}
\label{fig:15}
\end{minipage}\hfill
\begin{minipage}[t]{.48\textwidth}
\vspace{0pt}\centering
\includegraphics[width=\linewidth]{figures/fig16.pdf}
\caption{Association judgments for 32 inputs across four source chains, crossed by coherent or altered association, flat or linked presentation, and forward (F) or swapped (S) answer order. Letters retain the saved A/B responses; circles indicate correct associations and diamonds incorrect associations.}
\label{fig:16}
\end{minipage}
\end{figure*}

\subsubsection{Diagnostics turn evidence gaps into concrete research actions}

Broken evidence chains must be translated into targeted, physically grounded feedback to effectively drive autonomous corrective actions. When confronted with voluminous unstructured text outputs from quantum-chemical calculations, large language models readily overlook critical geometric and physical anomalies within their attention context. In contrast, structured diagnostic operators eliminate textual redundancy by isolating specific failure modes (e.g., steric clash atom pairs, imaginary-mode misalignments, or trajectory endpoint dislocations), providing concrete targets for subsequent local iterations.

Autonomous exploration trajectories highlight the important influence of diagnostic granularity on recovery. In a Brief-errors HA1 trial, where specific diagnostic feedback on IRC-endpoint linkage was ablated in favor of compact error codes, the Computational Agent succumbed to blind retries—mechanically substituting source-minimum references while leaving the underlying geometric disconnection unresolved. Conversely, in a Guided NB1 trial, the underlying verification detected zero imaginary frequencies and categorized the failure as a normal-mode error; the Coordinator immediately formulated a targeted retry plan, directing the Computational Agent to revise transition-state guess parameters and optimization constraints, which were faithfully incorporated into the rescheduled calculation. This case comparison can demonstrate that recovery in long-horizon mechanism elucidation depends not merely on detecting setbacks, but on bridging evidence defects directly to the specific chemical entities and computational parameters requiring revision.\citep{ref18}

This mechanism substantiates the empirical efficacy of Cloudra's multi-agent division of responsibilities: Computational Agents iteratively refine specific geometries and paths using localized feedback, the Coordinator oversees branching hypotheses and route selection across the reaction network, and the Supervisor audits cumulative support against immutable task rubrics. Supported by central asynchronous execution and durable ledgers, each role can maintain continuous context around shared chemical entities, reducing communicative drift across extended exploration horizons.

\begin{figure*}[tbp]
\centering
\includegraphics[width=.96\textwidth]{figures/fig17.pdf}
\caption{Completion requests and DAG feedback in the original IA1 trajectory. A report with unresolved evidence is followed by a completion request, a response identifying missing short\_irc proof, feedback to the model, and a revised partial-completion request.}
\label{fig:17}
\end{figure*}

\subsubsection{Outstanding requirements must remain part of completion decisions}

The final stage of scientific delivery requires a global determination: whether the accumulated computational artifacts genuinely satisfy rigorous physicochemical acceptance contracts. An autonomous IA1 trajectory reveals an inherent vulnerability of large language models during scientific research: premature completion hallucination. In this case, the reaction step lacked mandatory bidirectional short\_irc verification, with an empty waiver list; however, because the Supervisor had previously offered affirmative comments on intermediate optimization convergence, the model suffered cognitive drift and attempted to declare the exploration complete. In an unconstrained or purely conversational architecture, such unsupported mechanistic claims would escape scrutiny and produce false deliverables.

Here, the delivery graph's deterministic physical gate intervened decisively. The harness rejected the premature completion request and injected the unresolved short\_irc requirement as a hard constraint back into the context. Confronted with this unyielding physical reality, the model immediately rectified its claim in the subsequent submission, correctly downgrading its status to partial completion (Figure~\ref{fig:17}).

When physical state is explicitly recorded in an immutable ledger, the model reliably anchors its reasoning to this authoritative source. Across 12 judgments (six pairs with reversed answer-option order) evaluating IA1 under both minimal summary and extended 12-message history contexts, the model achieved 100\% accuracy in identifying unsatisfied scientific requirements across all state queries and reversed option orders. The delivery DAG thus provides an immutable coordinate system for subsequent reasoning: demarcating verified entities, pinpointing existing deficits, and determining whether to pursue additional calculations or submit partial deliverables reflecting physical reality.


\input{sections/03_benchmark_discussion}
